%\ifodd\value{page}\hbox{}\newpage\fi
\chapter*{Śrī Lalitā Sahasranāma — Śloka 16}
\addcontentsline{toc}{chapter}{Śloka 16 - Nomi 32,33}

\begin{sloka}
{\sanskritfont
\begin{tabular}{c}
अरुणारुणकौसुम्भवस्त्रभास्वत्कटीतटी ।\\
रत्नकिङ्किणिकारम्यरशनादामभूषिता ॥ १६॥
\end{tabular}

\vspace{1.5em}

{\middlesize \iastfont
aruṇāruṇa-kauśumbha-vastra-bhāsvat-kaṭī-taṭī |\\
ratna-kiṅkiṇikā-ramya-raśanā-dāma-bhūṣitā ||16||}

\vspace{1em}

{\middlesize
\anvaya{%
अरुणारुणकौसुम्भवस्त्रभास्वत्कटीतटी ।
रत्नकिङ्किणिकारम्यरशनादामभूषिता ।
}
}

\middlesize \iastfont \itmean{%
«Il cui fianco alla vita risplende per le vesti color rosso intenso di \textit{kauśumbha};  
ornata di cinture e catene, rese deliziose da campanellini preziosi.»%
}

}
\end{sloka}

\wordmeaning{%
\wordline{अरुण-अरुण-कौसुम्भ-वस्त्र-भास्वत्-कटी-तटी}
{aruṇa-aruṇa-kauśumbha-\\vastra-bhāsvat-kaṭī-taṭī}
{ il cui fianco della vita (\textit{kaṭī-taṭī}) è splendente (\textit{bhāsvat}) per le vesti (\textit{vastra}) tinte di \textit{kauśumbha}, di colore rosso intenso (\textit{aruṇa-aruṇa}); }%

\wordline{रत्न-किङ्किणिका-रम्य-रशना-दाम-भूषिता}
{ratna-kiṅkiṇikā-ramya-raśanā-\\dāma-bhūṣitā}
{ ornata (\textit{bhūṣitā}) di cinture (\textit{raśanā}) e catene (\textit{dāma}), rese deliziose (\textit{ramya}) da campanellini preziosi (\textit{ratna-kiṅkiṇikā}); }%
}

\textit{Śloka} 16 concentra la contemplazione sulla regione della vita e dei fianchi. Il composto \textit{aruṇāruṇa-kauśumbha-vastra} non descrive semplicemente un colore, ma una qualità di presenza: il rosso intenso delle vesti rituali rende la forma luminosa e immediatamente percepibile.

Nel secondo \textit{pāda}, l’ornamento della cintura è caratterizzato dalla grazia del dettaglio sonoro: le \textit{kiṅkiṇikā}, campanellini preziosi, rendono “ramya” l’intero insieme di cintura e catena. Il movimento e il tintinnio non aggiungono solo luce, ma ritmo.

Il centro del corpo è avvolto da una \textit{Śakti} rossa, cioè da un’energia già consacrata. Il confine tra alto e basso non è presentato come zona di colpa o di ambiguità, ma come una cintura di luce, luogo di passaggio e di misura. Il limite, rappresentato dalla cintura, non imprigiona né trattiene: dà forma. Il suo suono non è accessorio ornamentale, ma ritmo, rendendo percepibile il movimento della Dea come cadenza ordinata e sacra. In questo modo, la regione più delicata e culturalmente ambivalente del corpo femminile viene trasfigurata in una soglia luminosa e ritmica, attraverso la quale la potenza di \textit{Lalitā} passa, si contiene e si offre.

Questo \textit{śloka} contiene due \textit{nāma}.  
\textbf{\textit{Aruṇāruṇa-kauśumbha-vastra-bhāsvat-kaṭī-taṭī}} descrive la vita della Dea come risplendente del rosso rituale.  
\textbf{\textit{Ratna-kiṅkiṇikā-ramya-raśanā-dāma-bhūṣitā}} presenta la cintura come ornamento prezioso e sonoro, definito dalla grazia del suo movimento.
